\documentclass[10pt,UTF8]{beamer}

\usetheme{metropolis}
\usepackage{amsmath,amssymb,amsfonts}
\usepackage{booktabs}
\usepackage{graphicx}
\usepackage{xcolor}
\graphicspath{{./figures/}}

\title{Generative Modeling of\\[3pt] Real-Time Stochastic Dynamics}
\date{\today}
\author{Yang-yang Tan}

\definecolor{burgundy}{RGB}{127,0,32}
\setbeamercolor{alerted text}{fg=burgundy}
\setbeamercolor{block title}{use=normal text,fg=normal text.fg,bg=normal text.bg!80!fg}

\begin{document}
\maketitle

\begin{frame}{Dynamic critical behavior}
\textbf{Hohenberg--Halperin viewpoint}
\begin{align*}
  \partial_t\phi_a(\mathbf x,t)
  &=-\sum_b\mathcal K_{ab}[\phi]\,
  \frac{\delta F[\phi]}{\delta\phi_b(\mathbf x,t)}
  +\zeta_a(\mathbf x,t),\\
  \left\langle\zeta_a(\mathbf x,t)\zeta_b(\mathbf x',t')\right\rangle
  &=2T\mathcal K_{ab}\,
  \delta^{(d)}(\mathbf x-\mathbf x')\delta(t-t').
\end{align*}

\begin{block}{Model B with a continuum \(\phi^4\) free energy}
\begin{align*}
  F[\phi]
  &=\int d^2x\left[
    \frac{1}{2}(\boldsymbol\nabla\phi)^2
    +\frac{r}{2}\phi^2+\frac{u}{4}\phi^4\right],\\
  \mathcal K&=-\Gamma\nabla^2,
  \qquad
  \partial_t\phi
  =\Gamma\nabla^2\frac{\delta F}{\delta\phi}+\zeta,\\
  M(t)&=\frac{1}{|\Omega|}\int_\Omega\phi(\mathbf x,t)\,d^2x=M(0).
\end{align*}
\end{block}
\end{frame}

\begin{frame}{Learning a finite-time transition}
\begin{block}{Physical training pairs}
\begin{equation*}
  \phi\equiv\phi_t,
  \qquad
  \Delta\phi\equiv\phi_{t+\tau}-\phi_t,
  \qquad
  \tau=10 .
\end{equation*}
\end{block}

\begin{block}{Learned conditional transition}
\begin{equation*}
  p_\theta(\Delta\phi\mid\phi)
  \simeq
  p_\tau(\phi_{t+\tau}-\phi_t\mid\phi_t).
\end{equation*}
The condition \(\phi\) remains clean; the diffusion model generates only
\(\Delta\phi\).
\end{block}

\begin{block}{Trajectory generation}
\begin{align*}
  \Delta\phi_n\sim p_\theta(\Delta\phi\mid\phi_{n\tau}),
  &\qquad
  \phi_{(n+1)\tau}=\phi_{n\tau}+\Delta\phi_n,\\
  n=0,\ldots,19
  &\qquad\Longrightarrow\qquad
  t_{\max}=20\tau=200 .
\end{align*}
\end{block}
\end{frame}

\begin{frame}{Spatial evolution from a Gaussian initial field}
\centering
\vfill
\makebox[\textwidth][c]{%
  \includegraphics[width=1.12\textwidth]{gaussian_spatial_evolution.pdf}}
\vfill
{\scriptsize
The two rows start from the same configuration and are evolved with independent
stochastic noise.}
\vfill
\end{frame}

\begin{frame}{Equilibrium HMC initial fields}
\centering
\includegraphics[width=\textwidth,height=0.73\textheight,keepaspectratio]
  {hmc_observables.pdf}

\vspace{-0.1cm}
{\scriptsize
At \(t=200\): \(\phi_2\) \(+1.84\%\), \(\phi_4\) \(+2.88\%\), centered
variance \(-0.78\%\).  The direct dynamics stays stationary, while the DM
develops a smooth accumulated bias.}
\end{frame}

\begin{frame}{Nonequilibrium Gaussian initial fields}
\centering
\includegraphics[width=\textwidth,height=0.73\textheight,keepaspectratio]
  {gaussian_observables.pdf}

\vspace{-0.1cm}
{\scriptsize
The DM captures the initial decrease and subsequent growth of fluctuations.
At \(t=200\): \(\phi_2\) \(+3.88\%\), \(\phi_4\) \(+7.17\%\), centered
variance \(+3.81\%\).}
\end{frame}

\begin{frame}{Time correlations from HMC initial fields}
\centering
\includegraphics[width=\textwidth,height=0.73\textheight,keepaspectratio]
  {hmc_lag_average.pdf}

\vspace{-0.05cm}
{\scriptsize
The relaxation hierarchy is reproduced.  The main systematic difference is
in the slowest mode: at lag \(10\), \(\rho_1=0.8093(11)\) for DM and
\(0.7825(12)\) for Model B.}
\end{frame}

\begin{frame}{Time correlations from Gaussian initial fields}
\centering
\includegraphics[width=\textwidth,height=0.73\textheight,keepaspectratio]
  {gaussian_lag_average.pdf}

\vspace{-0.05cm}
{\scriptsize
The ordering of relaxation scales remains correct, but the DM retains excess
long-time memory in selected low modes.}
\end{frame}

\end{document}
